% !PRINTED 582
\weberpage{582}
On the question which of the three roots of the cubic
equation (19) is to be set for $\xi = 3\eta + 1$, which for $3\eta_1 + 1$ or for
$3\eta_2 + 1$, this much may be said in general: the sign of the
product (17) must agree with the sign of (18), which is determined
by the definition (1), (2), though it still depends on the choice
of the primitive root $g$. By this, three of the six possible
assignments are excluded. Which assignment among the three
remaining is to be made depends on the choice of $r$, and would
give occasion to investigations similar to those carried out in the
previous section for the determination of the sign of the square root\footnote{\gesperrt{Kummer}, Journ. f. Mathematik, Bd. 32.}.

We emphasise further the theorem proved by the formulae
(14), (16):

\gesperrt{If} $n$ \gesperrt{is a prime of the form} $3k + 1$, \gesperrt{then} $n$ \gesperrt{is}
\gesperrt{representable by the form} $a^2 - ab + b^2$ \gesperrt{and} $4n$ \gesperrt{by the form}
$A^2 + 27B^2$, \gesperrt{where} $a, b, A, B$ \gesperrt{are integers.}

From this it follows further that $n$ is also representable in
the form $x^2 + 3y^2$. For if of the two numbers $a, b$ one,
say $a$, is even, then
\[
n = \left(\frac{a}{2} - b\right)^2 + \tfrac{3}{4}a^2,
\]
and if $a$ and $b$ are both odd, then
\[
n = \tfrac{1}{4}(a + b)^2 + \tfrac{3}{4}(a - b)^3. % [sic] printed exponent 3, expected 2
\]

The two-term periods of the $9$th roots of unity also
satisfy a cubic equation, and because this equation plays an
important role in the general theory of the cubic cyclotomic
fields to be treated later, we shall consider it here for the
sake of completeness, although it properly belongs to a more
general region, in which the degrees of the roots of unity
are no longer primes.

A $9$th root of unity $r$ satisfies the equation of degree $6$ % [sic] the German prints "Einheitswuzel" for "Einheitswurzel"
\weq{21}{r^6 + r^3 + 1 = 0,}
and we have three conjugate two-term periods
\weq{22}{\eta = r + r^{-1}, \quad \eta' = r^2 + r^{-2}, \quad \eta'' = r^4 + r^{-4},}
which are the roots of the cubic equation
\weq{23}{\eta^3 - 3\eta + 1 = 0}
%CONT>>
